\bookchapter{Capítulo 17: Integración del móvil con otros dispositivos}

\chapterintro{El móvil en tu mano es en realidad un supermando a la vez que centro de datos. No solo hace llamadas: puede conectar el ordenador, la TV, el reloj, los altavoces y el coche, para que estos dispositivos dejen de ser islas. En este capítulo hablaremos de cómo hacerlos “trabajar en equipo”.}

\booksection{Sincronización con el ordenador: el móvil y el PC ya no son dos mundos separados}

Sinceramente, lo que más me molestaba era hacer una foto con el móvil y querer editarla en el ordenador: tenía que buscar el cable, esperar la transferencia. Luego descubrí que en realidad no hace falta cable; una buena integración es diez veces más rápida.

\begin{itemize}
  \item \textbf{Ecosistema Apple (AirDrop + Portapapeles universal)}:Si usas iPhone con Mac, AirDrop es lo primero que debes aprender. Toca compartir > AirDrop > elige tu Mac; la foto pasa al instante, sin datos ni Wi-Fi, usando Bluetooth y Wi-Fi directos. Hay otra función fácil de ignorar: "Portapapeles universal" —copias texto en el móvil y lo pegas directamente en el Mac, como magia.
  \item \textbf{Android + Windows (Conectar a Windows)}:Microsoft tiene una App llamada "Conectar a Windows", preinstalada en Samsung; para otros Android, descárgala de la tienda. Tras emparejar, puedes ver notificaciones del móvil, responder SMS, ver fotos e incluso recibir y hacer llamadas desde el ordenador. Llevo dos años usándola; en reuniones no necesito mirar el móvil constantemente: basta con la notificación en la pantalla del PC.
  \item \textbf{Colaboración multipantalla de marca}:La "colaboración multipantalla" de Huawei y "Xiaomi Magic Share" son aún más agresivas: proyectan la pantalla del móvil en el ordenador. Puedes arrastrar archivos con el ratón, usar Douyin o responder WeChat desde el PC. Copiar/pegar entre dispositivos, arrastrar archivos: todo funciona con una fluidez increíble.
\end{itemize}

\booksection{Proyectar a la TV: ver todos juntos}

En las fiencias de Año Nuevo, mi prima quería poner su video de viaje de graduación en la tele. Después de mucho intentar, resultó ser muy simple:

\begin{itemize}
  \item \textbf{iPhone TV}:Desliza desde la esquina superior derecha el centro de control, toca el icono de dos cuadrados apilados (Espejo de pantalla), encuentra tu TV y conéctate. Requisito: TV y móvil en la misma Wi-Fi, y la TV soporta AirPlay (la mayoría de smart TVs lo hacen).
  \item \textbf{Android TV}:Cada marca lo llama distinto: Huawei "proyección inalámbrica", Xiaomi "proyección", OPPO/vivo "interacción multipantalla". Normalmente está en el panel de accesos rápidos; si no, búscalo en ajustes.
  \item \textbf{Reproducción directa desde App de vídeo}:Si solo quieres ver iQiyi, Bilibili o YouTube en la tele, es más simple: en la interfaz de reproducción busca un pequeño icono rectangular con símbolo Wi-Fi, tócalo, elige la TV y el video se transfiere. El móvil queda libre para otras cosas mientras el video sigue en la tele.
  \item \textbf{¿Qué hacer si falla la proyección?}:El 90\% de los fallos es por no estar en la misma Wi-Fi. El otro 10\% es que la función de proyección de la TV no está activa; ve a los ajustes de la TV y busca "interacción multipantalla" o "AirPlay".
\end{itemize}

\booksection{Con smartwatch/pulsera inteligente: el gestor de salud en tu muñeca}

Le compré una pulsera inteligente a mi madre; al principio dijo “¿para qué si no hago ejercicio?”. Tras dos semanas de uso, descubrió que solo dormía cinco horas al día y empezó a ajustar su rutina. Ese es el valor más práctico del smartwatch/pulsera: convertir datos de salud invisibles en números visibles.

\begin{itemize}
  \item \textbf{Filtro de notificaciones}: Cuando el móvil vibra sin parar en el bolso, alza la muñeca y mira: ¿es el WeChat del jefe o una promoción de compras? Lo decides en un segundo; si no merece sacar el móvil, déjalo.
  \item \textbf{Sincronización de datos de salud}: El reloj monitoriza las 24 horas el ritmo cardíaco, el oxígeno en sangre, el sueño y los pasos, y los sincroniza automáticamente con la App de salud del móvil. Con el tiempo verás tus tendencias: si caminas poco este mes o si duermes mejor esta semana.
  \item \textbf{Sincronización de tarjetas de transporte/acceso}: La tarjeta de transporte y la de acceso vinculadas en el móvil se sincronizan solas con el reloj en la mayoría de marcas. En el metro no sacas el móvil: acercas la muñeca al torno y pasas. De verdad da la sensación de que «el futuro ya llegó».
  \item \textbf{"Buscar móvil" desde el reloj}: El reloj suele traer la función «buscar móvil»; tócala y el móvil sonará. Yo la uso mucho más de lo que esperaba: dejo el móvil en cualquier lado y se me olvida.
\end{itemize}

\booksection{Con hogar inteligente: el móvil es el “control maestro” de tu casa}

El Hogar inteligente suena misterioso, pero en realidad son tres cosas: control remoto desde el móvil, reglas de automatización y control por voz.

Te cuento algúnos usos prácticos propios y verás lo útil que es:

\begin{itemize}
  \item \textbf{Aire acondiciónado}:De vuelta a casa en verano, enciende el aire acondiciónado desde el móvil antes de llegar; entrarás a un ambiente fresco. Cuesta poco más de 100 yuanes comprar un enchufe inteligente para el aire; no hace falta cambiar el aparato.
  \item \textbf{Luces}: Acostado en la cama, sin ganas de levantarte a apagar la luz: toca el móvil o di «Hey Siri / Xiao Ai, apaga la luz».
  \item \textbf{robot aspiradora}:Antes de salir, toca "iniciar limpieza" en el móvil; al volver, el suelo estará limpio. También puedes programar limpieza diaria.
  \item \textbf{Cámaras}: Si hay ancianos, niños o mascotas en casa, puedes echar un vistazo al hogar desde el móvil en cualquier momento.
\end{itemize}

¿Qué ecosistema elegir? En resumen:

\begin{itemize}
  \item Usas iPhone prioriza dispositivos compatibles con Apple HomeKit (más caros pero mejor privacidad).
  \item Buscas relación calidad-precio Xiaomi Mijia tiene el ecosistema más completo: lámparas, enchufes, robots aspiradora, etc.
  \item Usas Huawei prioriza HarmonyOS Connect; la mejor experiencia con dispositivos Huawei.
\end{itemize}

\booksection{Auriculares y altavoces Bluetooth: la libertad del wireless}

Cambiar auriculares con cable por Bluetooth fue para mí la mejora más notable en la experiencia del móvil. Nunca más un ovillo de cables en el bolso ni preocuparse de que se enganchen en las manillas.

\begin{itemize}
  \item \textbf{Emparejar es sencillo}: Los auriculares entran solos en modo de emparejamiento la primera vez que se encienden (luz parpadeante); activas el Bluetooth del móvil, los detectas, tocas y conectan. Después se enlazan solos cada vez: no haces nada.
  \item \textbf{Misma marca: ventajas extra}: Los AirPods con iPhone muestran una ventana con la batería y cambian sin costuras entre dispositivos Apple. Los Huawei FreeBuds con Huawei y los auriculares Xiaomi con Xiaomi tienen ventajas similares.
  \item \textbf{Altavoz Bluetooth}: En reuniones familiares o excursiones, un altavoz Bluetooth duplica el ambiente. Conectado, tu móvil es tu mesa de DJ.
  \item \textbf{Ojo con la latencia}: Jugando con auriculares Bluetooth puede haber un ligero retraso (el disparo suena medio segundo tarde). Si eres sensible a la latencia, usa cable o compra unos Bluetooth con «modo de baja latencia».
\end{itemize}

\begin{tcbinfo}{Tecnologia UWB (ultra wideband): una "sensacion espacial" mas precisa que el Bluetooth}
Ademas del Bluetooth, los telefonos nuevos esconden un chip UWB (Ultra-Wideband). Su capacidad es la\textbf{medicion de distancia y posicionamiento espacial precisos}-- calcula a que distancia esta otro dispositivo y en que direccion, con precision centimetrica. El Bluetooth solo dice “cerca”; UWB dice “a 1,2 m a tu izquierda”. Aplicaciones:

\begin{itemize}
  \item \textbf{AirTag / busqueda precisa}: Al buscar un AirTag, el iPhone muestra una flecha con "7 m, gira a la izquierda" gracias a UWB. El Bluetooth solo indica si esta o no en el rango.
  \item \textbf{Llave digital del coche}: Los iPhone y gamas altas de Android con UWB funcionan como llave: al acercarse se desbloquea y al alejarse se bloquea; distingue la posicion del conductor y evita ataques de retransmision.
  \item \textbf{Navegacion en interiores}: En centros comerciales y aeropuertos, balizas UWB ofrecen navegacion a nivel de pasillo, mucho mas precisa que el GPS.
  \item \textbf{Modelos compatibles}: Desde iPhone 11 todos llevan chip UWB (U1/U2); en Android, algunas gamas altas de Xiaomi, Samsung, vivo y OPPO ya lo incorporan.
\end{itemize}

\end{tcbinfo}

\begin{tcbinfo}{LE Audio y Auracast: la nueva generacion de audio por Bluetooth}
El clasico audio Bluetooth (A2DP) usado durante veinte anos esta siendo sustituido por\textbf{LE Audio}(audio de bajo consumo), el nuevo estandar desde Bluetooth 5.2:

\begin{itemize}
  \item \textbf{Menor consumo y mejor sonido}: Usa el nuevo codec LC3; mejor sonido que el viejo SBC y menor consumo; la autonomia de los auriculares mejora notablemente.
  \item \textbf{Audio de doble via}: Los auriculares izquierdo y derecho se transmiten por separado; al cambiar a un solo lado no se corta.
  \item \textbf{Audio de difusion Auracast}: La funcion estrella de LE Audio. Un telefono o televisor puede emitir a\textbf{numerosos}auriculares a la vez: en aeropuertos, museos, gimnasios o aulas, un emisor y cada persona con sus propios auriculares escucha lo mismo.
  \item \textbf{Situación actual}: Los auriculares y telefonos de gama alta lanzados en 2024-2026 soportan LE Audio; los dispositivos antiguos necesitan actualizacion de hardware y no se desbloquea por software. Al comprar auriculares nuevos, comprueba si llevan la marca LE Audio/Auracast.
\end{itemize}

\end{tcbinfo}

\booksection{Con sistemas del coche: compañero de seguridad al conducir}

Usar el móvil mientras conduces es ilegal y peligroso. Pero a veces necesitas navegación, llamadas o música. Ahí entran CarPlay (Apple) y Android Auto.

\begin{itemize}
  \item \textbf{Cómo conectar}: Si tu coche lo soporta (la mayoría de los recientes), conecta el móvil al USB del coche con cable y la pantalla cambia sola a CarPlay o Android Auto. Algunos coches nuevos soportan conexión inalámbrica.
  \item \textbf{Qué puedes hacer}: Navegación (mucho mejor que la del coche), llamadas (contesta desde el volante), música (las Apps del móvil se proyectan al coche) y WeChat por voz (Siri o Xiao Ai leen y responden por ti).
  \item \textbf{Coches antiguas también pueden}: Si tu coche es viejo y no lo soporta, por unos cientos de yuanes compra una pantalla CarPlay externa, pégala en el salpicadero y la experiencia es parecida.
\end{itemize}

\booksection{Seguridad y privacidad en la conectividad}

Cuantos más dispositivos conectes, mayor es el riesgo de filtración de privacidad. Algunos puntos a tener en cuenta:

\begin{itemize}
  \item Antes de proyectar en público, revisa primero las notificaciones de tu móvil para que no haya mensajes privados; no dejes que toda la familia vea lo que no deben ver.
  \item Antes de regalar o vender el dispositivo, desemparejalo primero: ve a Ajustes del móvil, encuentra ese dispositivo y selecciona "Eliminar" o "Desparejar".
  \item Para cámaras de hogar inteligente, compra de marca con cifrado; evita las genericas baratas. He oído demasiadas noticias de cámaras hackeadas con transmisiones en vivo.
\end{itemize}

\begin{tcbpractice}{Practica: proyecta la pantalla del móvil a la TV y configura una escena inteligente}
\begin{enumerate}
  \item Asegúrate de que el móvil y la TV están en la misma Wi-Fi.
  \item En la TV, busca y activa "Espejo de pantalla" o "AirPlay".
  \item En el móvil: iPhone, centro de control > Espejo de pantalla. Android, desliza y busca "proyección inalámbrica".
  \item Elige tu foto favorita y proyéctala en la TV para que tu familia la vea.
  \item si tienes dispositivos de Hogar inteligente, intenta crear un "modo fuera de casa" en la App del móvil: un toque para apagar luces, apagar aire aencionado, iniciar robot aspiradora.
\end{enumerate}

\end{tcbpractice}

\begin{tcbinfo}{Caso 1: optimizacion de autonomia en la practica}
\begin{itemize}
  \item \textbf{Escenario}: Un usuario de Xiaomi 14 tiene mala autonomia y se queda sin bateria a mediodia.
  \item \textbf{Pasos}: 1) Ajustes > Bateria > ranking de consumo; 2) identifica apps anomalias en segundo plano; 3) cambia de 5G a 4G; 4) restringe segundo plano; 5) activa modo de ahorro; 6) al cargar, quita la funda para disipar calor.
  \item \textbf{Resultado esperado}: La autonomia se prolonga un 30-50\%; se puede leer el estado de salud de la bateria.
\end{itemize}

\end{tcbinfo}
